\documentclass[aspectratio=169,12pt]{beamer}

\usepackage[T1]{fontenc}
\usepackage[utf8]{inputenc}
\usepackage{lmodern}
\usepackage{amsmath,amssymb,bm}
\usepackage{graphicx}
\usepackage{booktabs,array,multicol}
\usepackage{hyperref}
\usepackage{xcolor}
\usepackage{ragged2e}
\usepackage{microtype}

% --------------------------------------------------
% Academic beamer style for a broad-audience lecture
% Design principle: every important concept has two layers:
% 1) Core intuition for first-time AI learners
% 2) Technical lens for students with stronger background
% --------------------------------------------------
\definecolor{LectureBlue}{RGB}{51,51,180}
\definecolor{FooterBlue}{RGB}{24,26,96}
\definecolor{DarkText}{RGB}{20,20,20}
\definecolor{LightBox}{RGB}{246,246,250}
\definecolor{SoftBlue}{RGB}{236,240,255}
\definecolor{SoftGreen}{RGB}{238,248,240}
\definecolor{SoftOrange}{RGB}{255,246,235}
\definecolor{MidGray}{RGB}{85,85,85}
\definecolor{DeepGreen}{RGB}{20,96,64}
\definecolor{DeepOrange}{RGB}{150,82,20}

\setbeamertemplate{navigation symbols}{}
\setbeamertemplate{headline}{%
  \leavevmode%
  \hbox{%
    \begin{beamercolorbox}[wd=\paperwidth,ht=0.95cm,dp=0.20cm,leftskip=0.32cm]{headline}%
      \usebeamerfont{headline}\color{white}\currentframetitle
    \end{beamercolorbox}%
  }%
}
\setbeamercolor{headline}{bg=LectureBlue,fg=white}
\setbeamerfont{headline}{series=\mdseries,size=\Large}
\setbeamertemplate{frametitle}{}

\setbeamertemplate{footline}{%
  \leavevmode%
  \hbox{%
    \begin{beamercolorbox}[wd=.78\paperwidth,ht=0.32cm,dp=0.2cm,leftskip=0.22cm]{footerleft}%
      \color{white}\scriptsize don-huang.com
    \end{beamercolorbox}%
    \begin{beamercolorbox}[wd=.22\paperwidth,ht=0.32cm,dp=0.2cm,rightskip=0.22cm plus1fil]{footerright}%
      \color{white}\scriptsize \hfill \insertframenumber\ / \inserttotalframenumber
    \end{beamercolorbox}%
  }%
}
\setbeamercolor{footerleft}{bg=FooterBlue,fg=white}
\setbeamercolor{footerright}{bg=LectureBlue,fg=white}

\setbeamercolor{normal text}{fg=DarkText,bg=white}
\setbeamercolor{item}{fg=FooterBlue}
\setbeamertemplate{itemize item}{\color{FooterBlue}$\bullet$}
\setbeamertemplate{itemize subitem}{\color{FooterBlue}$\circ$}
\setbeamertemplate{enumerate item}{\color{FooterBlue}\insertenumlabel.}

\newcommand{\currentframetitle}{}
\newcommand{\slideTitle}[1]{\gdef\currentframetitle{#1}\frametitle{#1}}
\newcommand{\keyterm}[1]{\textbf{#1}}
\newcommand{\twocol}[2]{\begin{columns}[T,onlytextwidth]\begin{column}{0.49\textwidth}#1\end{column}\begin{column}{0.49\textwidth}#2\end{column}\end{columns}}
\newcommand{\mainidea}[1]{%
  \vspace{0.08cm}
  \begin{beamercolorbox}[wd=\linewidth,rounded=true,sep=0.18cm]{mainideabox}
    \small\textbf{Core intuition:} #1
  \end{beamercolorbox}%
}
\newcommand{\techlens}[1]{%
  \vspace{0.08cm}
  \begin{beamercolorbox}[wd=\linewidth,rounded=true,sep=0.18cm]{techlensbox}
    \small\textbf{Technical lens:} #1
  \end{beamercolorbox}%
}
\newcommand{\finlens}[1]{%
  \vspace{0.08cm}
  \begin{beamercolorbox}[wd=\linewidth,rounded=true,sep=0.18cm]{finlensbox}
    \small\textbf{Finance lens:} #1
  \end{beamercolorbox}%
}
\newcommand{\takeaway}[1]{%
  \vspace{0.08cm}
  \begin{beamercolorbox}[wd=\linewidth,rounded=true,sep=0.18cm]{takeawaybox}
    \small\textbf{Takeaway:} #1
  \end{beamercolorbox}%
}
\newcommand{\visual}[1]{%
  \vspace{0.12cm}
  \begin{beamercolorbox}[wd=\linewidth,rounded=true,sep=0.16cm]{visualbox}
    \scriptsize\textbf{Suggested visual:} #1
  \end{beamercolorbox}%
}

\newcommand{\sourcefigure}[3]{%
  \vspace{0.08cm}
  \begin{center}
    \includegraphics[width=#2\linewidth,height=0.62\paperheight,keepaspectratio]{#1}\\[-0.02cm]
    {\tiny\color{MidGray}Source: #3}
  \end{center}%
}

\setbeamercolor{mainideabox}{bg=SoftBlue,fg=DarkText}
\setbeamercolor{techlensbox}{bg=SoftGreen,fg=DeepGreen}
\setbeamercolor{finlensbox}{bg=SoftOrange,fg=DeepOrange}
\setbeamercolor{takeawaybox}{bg=LightBox,fg=DarkText}
\setbeamercolor{visualbox}{bg=LightBox,fg=MidGray}

\title{AI Application in FinTech}
\subtitle{From Machine Learning Foundations to Agentic WealthTech}
\author{Don HUANG Yao Tung, PhD \\ AQUMON}
\institute{HKUST Lecture Series}
\date{2026 Edition}

\begin{document}

% --------------------------------------------------
\begin{frame}
\slideTitle{AI Application in FinTech}
\vspace{0.65cm}
{\Huge \textbf{AI Application in FinTech}}\\[0.25cm]
{\Large From Machine Learning Foundations to Agentic WealthTech}\\[0.45cm]
{\large Don HUANG Yao Tung, PhD}\\[0.10cm]
{\large Chief Scientist and Co-founder at AQUMON}
\vfill
\end{frame}


% --------------------------------------------------
\section{Part I: AI Technology Foundations}



\begin{frame}
\slideTitle{The Original Question: Can Machines Think?}
\begin{itemize}
  \item Alan Turing reframed a philosophical question into an operational test.
  \item Instead of asking what thinking ``really is,'' he asked whether a machine can behave intelligently in conversation.
  \item Today, AI is broader than imitation. We care about systems that can perceive, predict, reason, plan, act and learn from feedback.
\end{itemize}
\mainidea{AI is not one technology. It is a family of methods for turning information into useful behavior.}
\end{frame}

\begin{frame}
\slideTitle{Visual Example: The Turing Test}
\sourcefigure{figures_pdf/turing_test_diagram.pdf}{0.68}{Juan Alberto Sánchez Margallo, Wikimedia Commons, CC BY 2.5.}
\vspace{-0.04cm}
{\scriptsize\textit{Lecture point:} Turing reframed intelligence as observable behavior in a controlled interaction.}
\end{frame}


\begin{frame}
\slideTitle{A Practical Definition of AI}
\begin{itemize}
  \item \keyterm{Artificial intelligence} refers to computational systems that perform tasks normally associated with human cognition.
  \item Examples: perception, language understanding, pattern recognition, decision-making and problem solving.
  \item In business, a practical definition is simpler: AI helps convert data and context into better decisions or actions.
\end{itemize}
\takeaway{The central question is not ``Is it intelligent?'' but ``Which decision process does it improve?''}
\end{frame}

\begin{frame}
\slideTitle{AI, Machine Learning, Deep Learning and Foundation Models}
\textbf{Plain-language view}
\begin{itemize}
  \item AI: the big umbrella.
  \item Machine learning: learn patterns from data.
  \item Deep learning: neural networks with many layers.
  \item Foundation models: large general models adapted to many tasks.
\end{itemize}


\end{frame}

\begin{frame}
\slideTitle{Visual Example: AI, ML and Deep Learning Relationship}
\sourcefigure{figures_pdf/ai_ml_dl.pdf}{0.53}{AI-ML-DL.svg, Wikimedia Commons, CC BY-SA 4.0.}
\vspace{-0.45cm}
{\scriptsize\textit{Lecture point:} This hierarchy helps beginners avoid mixing up the umbrella field, learning methods and deep neural networks.}
\end{frame}


\begin{frame}
\slideTitle{Five Abilities of Modern AI}
\begin{center}
\begin{tabular}{lll}
\toprule
Ability & Simple meaning & Finance example \\
\midrule
Perceive & read inputs & parse PDFs and tables \\
Predict & estimate future/label & default probability \\
Generate & create content & draft client report \\
Reason & connect steps & explain portfolio risk \\
Act & use tools/workflows & rebalance proposal \\
\bottomrule
\end{tabular}
\end{center}
\mainidea{Modern AI is powerful because these abilities can be combined in one workflow.}
\end{frame}

\begin{frame}
\slideTitle{Machine Learning as Learning from Examples}
\mainidea{Machine learning is like teaching by examples instead of writing every rule by hand.}
\begin{itemize}
  \item Traditional programming: human writes rules, computer follows rules.
  \item Machine learning: human provides examples and objectives, computer learns patterns.
  \item Example: instead of manually listing every fraud pattern, we train on historical transactions.
\end{itemize}
\finlens{This is useful in finance because many patterns are too complex or too fast-changing for manual rules alone.}
\end{frame}

\begin{frame}
\slideTitle{Machine Learning as Function Approximation}
\begin{itemize}
  \item A compact mathematical view:
\end{itemize}
\[
    f_\theta(x) \approx y.
\]
\begin{itemize}
  \item Input $x$: image, text, customer profile, transaction, market features or portfolio state.
  \item Output $y$: label, forecast, recommendation, action or generated content.
  \item Learning means choosing parameters $\theta$ so that $f_\theta$ works well on unseen data.
\end{itemize}
\techlens{Most practical ML problems are about selecting the right inputs, objective function, validation design and deployment controls.}
\end{frame}



\begin{frame}
\slideTitle{Traditional System vs ML}
\centering
\includegraphics[width=0.93\textwidth]{figures_pdf/Traditional_vs_ML.pdf}
\end{frame}

\begin{frame}
\slideTitle{Data, Objective, Optimization, Validation}
\begin{itemize}
  \item A model is not trained from data alone.
  \item Four ingredients matter:
  \begin{enumerate}
    \item \keyterm{Data}: what the model can learn from.
    \item \keyterm{Objective}: what the model is rewarded for doing.
    \item \keyterm{Optimization}: how the model parameters are fitted.
    \item \keyterm{Validation}: how we know it works outside training data.
  \end{enumerate}
\end{itemize}
\[
\text{Data} + \text{Objective} + \text{Optimization} + \text{Validation} \Rightarrow \text{Useful model}.
\]
\end{frame}

\begin{frame}
\slideTitle{Three Main Learning Paradigms}
\begin{center}
\begin{tabular}{llll}
\toprule
Paradigm & Everyday analogy & Typical output & Finance example \\
\midrule
Supervised & answer key & prediction & default probability \\
Unsupervised & grouping & clusters/factors & client segmentation \\
Reinforcement & trial and feedback & policy & execution strategy \\
\bottomrule
\end{tabular}
\end{center}
\mainidea{The paradigm depends on what kind of feedback the system receives: labels, structure, or rewards.}
\end{frame}

\begin{frame}
\slideTitle{Supervised Learning: Learning with an Answer Key}
\begin{itemize}
  \item Training data contains both inputs and known answers:
\end{itemize}
\[
\mathcal{D}=\{(x_i,y_i)\}_{i=1}^n.
\]
\begin{itemize}
  \item Regression predicts a number: return, volatility, loss, revenue.
  \item Classification predicts a category: fraud/not fraud, churn/not churn, approve/reject.
  \item The model is useful only if it generalizes to new examples.
\end{itemize}
\finlens{Credit scoring, fraud detection and customer churn prediction are classic supervised-learning problems.}
\end{frame}

\begin{frame}
\slideTitle{Supervised Learning Examples}
\scriptsize
\begin{center}
\begin{tabular}{p{0.23\linewidth}p{0.29\linewidth}p{0.30\linewidth}p{0.11\linewidth}}
\toprule
Task & Input $x$ & Target $y$ & Output type \\
\midrule
Fraud detection & transaction amount, merchant, device, location & fraud label & probability \\
Credit risk & borrower profile, income, repayment history & default event & probability \\
Equity factor model & valuation, quality, momentum features & future return rank & score \\
Client churn & activity, holdings, service history & churn/not churn & class \\
\bottomrule
\end{tabular}
\end{center}
\takeaway{The same learning framework appears in very different financial products.}
\end{frame}

\begin{frame}
\slideTitle{Loss Function and Empirical Risk}
\begin{itemize}
  \item A model is trained by minimizing average loss on training data:
\end{itemize}
\[
\hat{\theta}=\arg\min_\theta \frac{1}{n}\sum_{i=1}^{n} L(f_\theta(x_i),y_i)+\lambda\Omega(\theta).
\]
\begin{itemize}
  \item $L$ measures prediction error.
  \item $\Omega$ regularizes model complexity.
  \item $\lambda$ controls the trade-off between fit and simplicity.
\end{itemize}
\techlens{Changing the loss function changes what the model optimizes. In finance, optimizing the wrong proxy can create the wrong behavior.}
\end{frame}

\begin{frame}
\slideTitle{Overfitting: The Most Important Warning}
\mainidea{Overfitting means the model memorizes noise instead of learning a repeatable pattern.}
\begin{itemize}
  \item A model can look excellent on historical data and fail in the future.
  \item Finance is especially vulnerable because market signals are weak and noisy.
  \item More complexity is not always better; sometimes it simply gives noise a better suit.
\end{itemize}
\[
\mathbb{E}[(\hat{f}(x)-y)^2] = \text{Bias}^2 + \text{Variance} + \text{Noise}.
\]
\takeaway{The real exam is not training performance; it is out-of-sample performance.}
\end{frame}

\begin{frame}
\slideTitle{Visual Example: Underfitting vs. Overfitting}
\sourcefigure{figures_pdf/bias_variance.pdf}{0.96}{scikit-learn documentation, Underfitting vs. Overfitting example, BSD-3-Clause.}
\vspace{-0.04cm}
{\scriptsize\textit{Lecture point:} The best model is not the most complex one; it is the one that generalizes.}
\end{frame}


\begin{frame}
\slideTitle{Train, Validation and Test Sets}
\begin{itemize}
  \item \keyterm{Training set}: fit the model.
  \item \keyterm{Validation set}: tune hyperparameters and select models.
  \item \keyterm{Test set}: estimate final out-of-sample performance.
  \item In time-series finance, random splitting can leak future information.
  \item Use chronological split, walk-forward validation or purged cross-validation.
\end{itemize}
\end{frame}
% --------------------------------------------------




\begin{frame}
\slideTitle{Unsupervised Learning: Finding Structure Without Labels}
\mainidea{Unsupervised learning is like organizing a messy library without a catalog. The system looks for natural structure.}
\begin{itemize}
  \item Data has no labels:
\end{itemize}
\[
\mathcal{D}=\{x_i\}_{i=1}^{n}.
\]
\begin{itemize}
  \item Common methods: K-means, hierarchical clustering, PCA and autoencoders.
  \item Finance examples: market regimes, client segmentation, peer grouping and style factor extraction.
\end{itemize}
\end{frame}

\begin{frame}
\slideTitle{K-Means Clustering}
\begin{itemize}
  \item K-means partitions observations into $K$ clusters.
  \item It minimizes within-cluster squared distance:
\end{itemize}
\[
\min_{C_1,\ldots,C_K}\sum_{k=1}^K \sum_{x_i\in C_k} \|x_i-\mu_k\|^2.
\]
\begin{itemize}
  \item Plain example: group customers by investment behavior.
  \item Finance example: cluster stocks by return correlation, valuation profile or risk behavior.
\end{itemize}
\techlens{Clusters are descriptions, not truth. They can change when regimes change or when the feature design changes.}
\end{frame}

\begin{frame}
\slideTitle{Visual Example: K-Means on Handwritten Digits}
\sourcefigure{figures_pdf/kmeans_digits.pdf}{0.72}{scikit-learn documentation, K-Means clustering on digits, BSD-3-Clause.}
\vspace{-0.04cm}
{\scriptsize\textit{Lecture point:} Clustering creates useful structure, but the structure depends on the representation and distance metric.}
\end{frame}
% --------------------------------------------------





\begin{frame}
\slideTitle{Reinforcement Learning: Learning by Feedback}
\mainidea{Reinforcement learning is learning by trying actions and receiving rewards. It is closer to training a player than fitting a spreadsheet formula.}
\begin{itemize}
  \item An agent observes state $s_t$, takes action $a_t$ and receives reward $r_t$.
  \item Objective: learn a policy $\pi(a|s)$ that maximizes cumulative reward.
\end{itemize}
\[
\max_\pi \; \mathbb{E}_\pi\left[\sum_{t=0}^{T}\gamma^t r_t\right].
\]
\finlens{Finance examples include order execution, market making, dynamic allocation and hedging.}
\end{frame}


\begin{frame}
\slideTitle{Why Reinforcement Learning Is Hard in Finance}
\begin{itemize}
  \item Feedback is noisy: good decisions can have bad short-term outcomes.
  \item Markets are non-stationary: yesterday's policy may fail tomorrow.
  \item Exploration is expensive: real trading mistakes cost money.
  \item The agent may affect the environment through its own trades.
\end{itemize}
\takeaway{RL is conceptually attractive in finance, but deployment needs simulation, constraints and human oversight.}
\end{frame}
% --------------------------------------------------
\begin{frame}
\slideTitle{Demo: Reinforcement Learning}
\footnotesize
\begin{beamercolorbox}[wd=\linewidth,rounded=true,sep=0.12cm]{visualbox}\scriptsize\textbf{Youtube: https://www.youtube.com/watch?v=Lt-KLtkDlh8} \end{beamercolorbox}
\centering
\end{frame}


\begin{frame}
\slideTitle{Neural Networks}
\mainidea{A neural network is a flexible function made by stacking many simple transformations.}
\begin{itemize}
  \item Each layer transforms the previous representation:
\end{itemize}
\[
    h^{(l)}=\sigma(W^{(l)}h^{(l-1)}+b^{(l)}).
\]
\begin{itemize}
  \item Early layers learn basic patterns; later layers learn task-specific representations.
  \item Deep learning became powerful because of data, GPUs, better optimization and open-source tooling.
\end{itemize}
\finlens{In finance, deep models can help with text, documents and high-dimensional data, but weak market signals still require careful validation.}
\end{frame}

\begin{frame}
\slideTitle{Visual Example: A Simple Neural Network Topology}
\sourcefigure{figures_pdf/artificial_neural_network.pdf}{0.66}{Cburnett, Wikimedia Commons, CC BY-SA 3.0 or GFDL.}
\vspace{-0.04cm}
{\scriptsize\textit{Lecture point:} A neural network stacks simple units into layers so the whole system can approximate complex functions.}
\end{frame}


\begin{frame}
\slideTitle{CNN, RNN and LSTM: Before Transformers}
\begin{itemize}
  \item \keyterm{CNN}: captures local spatial patterns; originally dominant in image tasks.
  \item \keyterm{RNN}: processes sequences step by step; useful for language and time series.
  \item \keyterm{LSTM}: improves long-range sequence learning through gated memory.
  \item These architectures prepared the ground for modern sequence models.
\end{itemize}
\takeaway{Architecture matters because it determines what patterns the model can learn efficiently.}
\end{frame}

\begin{frame}
\slideTitle{Attention Mechanism: The Key Intuition}
\mainidea{Attention lets the model decide which parts of the input are relevant to each other.}
\begin{itemize}
  \item In a sentence, a word may need information from another word far away.
  \item In a report, a conclusion may depend on numbers in a table several pages earlier.
  \item Attention creates a dynamic relevance map across tokens.
\end{itemize}
\finlens{For financial documents, attention helps connect terms, definitions, tables, footnotes and risk disclosures.}
\end{frame}



\begin{frame}
\slideTitle{Transformers}
\begin{itemize}
  \item Transformers are built from self-attention, feed-forward layers and normalization.
  \item They support parallel training over long sequences.
  \item They became the dominant architecture for language and increasingly for vision, audio, code and multimodal reasoning.
  \item Most modern LLMs are transformer-based or transformer-derived systems.
\end{itemize}
\visual{A clean transformer block illustration with token embeddings, self-attention and feed-forward layers.}
\end{frame}

\begin{frame}
\slideTitle{Visual Example: Transformer Architecture}
\sourcefigure{figures_pdf/transformer_architecture.pdf}{0.52}{Vaswani et al. (2017), Attention Is All You Need, NeurIPS paper Figure 1.}
\vspace{-0.04cm}
{\scriptsize\textit{Lecture point:} Transformers combine attention, feed-forward layers, residual connections and positional information.}
\end{frame}
% --------------------------------------------------




\begin{frame}
\slideTitle{Large Language Models}
\mainidea{An LLM learns language by predicting the next token at massive scale. Simple training objective; surprisingly rich behavior.}
\begin{itemize}
  \item Pretraining objective:
\end{itemize}
\[
\max_\theta \sum_t \log p_\theta(x_t|x_{<t}).
\]
\begin{itemize}
  \item Pretraining produces broad language and knowledge representations.
  \item Instruction tuning makes the model more helpful.
  \item Tool use and retrieval connect the model to external systems and fresh data.
\end{itemize}
\end{frame}




\begin{frame}
\slideTitle{Foundation Models: The New AI Layer}
\begin{itemize}
  \item Foundation models are trained broadly and adapted to many downstream tasks.
  \item They shift development from task-specific modeling to prompt, retrieval, tool and workflow design.
  \item They can support language, code, image, audio, video and structured data.
  \item In enterprise finance, the key architecture is often a controlled workflow around a foundation model, not the raw model alone.
\end{itemize}
\takeaway{Foundation models are general engines; financial value comes from connecting them to trusted data and responsible processes.}
\end{frame}

\begin{frame}
\slideTitle{Multimodal AI}
\begin{itemize}
  \item Multimodal models process combinations of text, image, audio, video and structured data.
  \item Finance examples:
  \begin{itemize}
    \item read PDFs, charts and tables;
    \item extract data from factsheets and reports;
    \item analyze earnings-call audio and transcripts;
    \item support advisor-client conversations with document grounding.
  \end{itemize}
\end{itemize}
\visual{Different inputs such as PDF, chart, table, voice and news flowing into a unified model, then producing a financial decision-support output.}
\end{frame}




\begin{frame}
\slideTitle{RAG: Retrieval-Augmented Generation}
\mainidea{RAG gives the model a trusted library before asking it to answer.}
\begin{enumerate}
  \item Parse and index documents.
  \item Retrieve relevant chunks.
  \item Generate answer with source references.
  \item Evaluate faithfulness and completeness.
\end{enumerate}
\finlens{RAG is central in finance because many answers must be traceable to policies, reports, contracts, fund documents or regulations.}
\end{frame}

\begin{frame}
\slideTitle{Visual Example: RAG Architecture}
\sourcefigure{figures_pdf/rag_architecture.pdf}{0.88}{Fan et al. (2024), \textit{A Survey on RAG Meeting LLMs: Towards Retrieval-Augmented Large Language Models}, Figure 3.}
\end{frame}


\begin{frame}
\slideTitle{RAG: What Can Go Wrong?}
\begin{itemize}
  \item Retrieval failure: the right document was not found.
  \item Chunking failure: the answer needs context split across chunks.
  \item Generation failure: the model misreads or over-interprets retrieved text.
  \item Evaluation failure: the answer sounds plausible but is not fully supported.
\end{itemize}
\techlens{A production RAG system needs document lineage, permissions, citation checks, evaluation sets and monitoring.}
\end{frame}
% --------------------------------------------------



\begin{frame}
\slideTitle{AI Agents}
\mainidea{An agent is not just a model that answers; it can plan, use tools, observe results and revise.}
\[
\text{Goal} \rightarrow \text{Plan} \rightarrow \text{Tool Use} \rightarrow \text{Observation} \rightarrow \text{Revision}.
\]
\begin{itemize}
  \item Agents are useful when a task requires multi-step execution rather than a single response.
  \item Examples: research assistant, coding assistant, portfolio diagnostic assistant.
  \item Finance needs permissioning, logging, sandboxing and approval gates.
\end{itemize}
\end{frame}




\begin{frame}
\slideTitle{World Models}
\mainidea{A world model learns how an environment evolves, so an agent can imagine consequences before acting.}
\begin{itemize}
  \item It supports counterfactual reasoning: ``What may happen if I take this action?''
  \item Compared with a text-only LLM, a world model focuses more on state, dynamics, action and simulation.
  \item It is especially important for agents that need to plan in changing environments.
\end{itemize}
\visual{An agent trying actions in a simulated environment before taking real-world actions.}
\end{frame}



\begin{frame}
\slideTitle{World Models: Mathematical View}
\begin{itemize}
  \item State: $s_t$; action: $a_t$; observation: $o_t$; reward: $r_t$.
  \item A world model approximates transition and observation structure:
\end{itemize}
\[
    \hat{p}(s_{t+1}, o_{t+1}, r_t \mid s_t, a_t).
\]
\begin{itemize}
  \item Planning can then be performed in imagination:
\end{itemize}
\[
    \max_{a_{0:T}} \; \mathbb{E}_{\hat{p}}\left[\sum_{t=0}^T \gamma^t r_t\right].
\]
\end{frame}

\begin{frame}
\slideTitle{World Models in Finance: Analogy and Limits}
\begin{itemize}
  \item Finance has an implicit world-model problem: investors reason about regimes and consequences.
  \item Potential uses:
  \begin{itemize}
    \item portfolio stress testing;
    \item scenario generation;
    \item execution simulation;
    \item synthetic market environments for strategy testing.
  \end{itemize}
  \item Major limitation: financial markets are adaptive systems; the model can change the world it tries to model.
\end{itemize}
\takeaway{A financial world model is useful only if it is treated as a scenario engine, not a crystal ball.}
\end{frame}
% --------------------------------------------------


\begin{frame}
\slideTitle{Evaluation of Modern AI Systems}
\begin{itemize}
  \item Classic ML evaluation: accuracy, AUC, RMSE, Sharpe ratio and cross-validation.
  \item Generative AI evaluation also needs:
  \begin{itemize}
    \item factuality and citation faithfulness;
    \item reasoning correctness;
    \item tool-use success rate;
    \item latency and cost;
    \item safety and compliance.
  \end{itemize}
\end{itemize}
\finlens{For finance, evaluation should be linked to the business decision being supported.}
\end{frame}

\begin{frame}
\slideTitle{AI Governance in Finance}
\begin{itemize}
  \item Financial AI systems need more than model accuracy.
  \item Key controls:
  \begin{itemize}
    \item data lineage and access control;
    \item model validation and monitoring;
    \item explainability and audit trail;
    \item human approval for regulated or high-impact decisions;
    \item cybersecurity and privacy protection.
  \end{itemize}
\end{itemize}
\takeaway{A good AI system should fail safely and leave evidence.}
\end{frame}

% --------------------------------------------------
\section{Part II: AI Applications in FinTech}

\begin{frame}
\slideTitle{Part II: AI Applications in FinTech}
\vspace{1.0cm}
\begin{center}
{\Huge AI Applications in FinTech}\\[0.35cm]
{\Large From automation to investment decision intelligence}
\end{center}
\end{frame}

\begin{frame}
\slideTitle{FinTech: Where AI Creates Value}
\begin{itemize}
  \item \keyterm{Scale}: serve more clients at lower marginal cost.
  \item \keyterm{Speed}: automate document reading, monitoring and reporting.
  \item \keyterm{Consistency}: reduce manual variation in decisions.
  \item \keyterm{Personalization}: adapt products and advice to client constraints.
  \item \keyterm{Insight}: extract signals from unstructured and high-dimensional data.
\end{itemize}
\mainidea{AI matters when it improves a workflow, not when it is merely attached to a product name.}
\end{frame}

\begin{frame}
\slideTitle{Major AI Use Cases in Financial Services}
\scriptsize
\begin{center}
\begin{tabular}{p{0.23\linewidth}p{0.31\linewidth}p{0.32\linewidth}}
\toprule
Area & AI task & Example output \\
\midrule
Customer service & summarization and dialogue & advisor copilot \\
Risk and fraud & anomaly detection & suspicious transaction alert \\
Credit & supervised prediction & default probability \\
Compliance & document review & policy breach flag \\
Research & retrieval and coding & factor test memo \\
WealthTech & optimization and personalization & portfolio proposal \\
Trading & sequential decision & execution schedule \\
\bottomrule
\end{tabular}
\end{center}
\end{frame}

\begin{frame}
\slideTitle{Why Finance Is Different}
\begin{itemize}
  \item Financial decisions are high-stakes and regulated.
  \item Data is noisy, non-stationary and often proprietary.
  \item Backtests can be misleading without strong controls.
  \item Model errors can create financial loss, conduct risk or regulatory risk.
  \item Therefore the AI workflow must combine automation with governance.
\end{itemize}
\takeaway{In finance, ``move fast and break things'' is not a strategy. It is an incident report waiting for a date.}
\end{frame}

\begin{frame}
\slideTitle{From Digital Finance to AI-Native Finance}
\begin{itemize}
  \item \keyterm{Digital finance}: online forms, dashboards, APIs and rule-based automation.
  \item \keyterm{AI-assisted finance}: models support search, summarization, prediction and monitoring.
  \item \keyterm{AI-native finance}: agents coordinate data, models, tools and approvals across an end-to-end process.
\end{itemize}
\mainidea{The next frontier is not just a chatbot. It is workflow intelligence.}
\end{frame}

\begin{frame}
\slideTitle{Robo-Advisory: Core Idea}
\mainidea{Robo-advisory converts client information into investment recommendations through a structured, repeatable process.}
\begin{enumerate}
  \item Client onboarding and risk profiling.
  \item Goal and constraint collection.
  \item Strategic asset allocation.
  \item Portfolio implementation.
  \item Monitoring, rebalancing and reporting.
\end{enumerate}
\finlens{The value comes from systematic advice, not from pretending human judgment is unnecessary.}
\end{frame}

\begin{frame}
\slideTitle{Visual Example: AI-Driven Robo-Advisory Roadmap}
\sourcefigure{figures_pdf/robo_advisor_roadmap.pdf}{0.88}{Feng, Li and Liu (2025), \textit{Robo-Advisors Beyond Automation}, Figure 2.}
\end{frame}


\begin{frame}
\slideTitle{Risk Profiling}
\begin{itemize}
  \item Risk profiling estimates a client's ability and willingness to take risk.
  \item Inputs may include age, income, wealth, horizon, liquidity needs, investment knowledge and loss tolerance.
  \item Common output: a risk score or risk category.
  \item AI can assist by detecting inconsistent answers, explaining trade-offs and mapping goals to portfolios.
\end{itemize}
\techlens{Risk profiling is partly statistical and partly behavioral. The hardest part is not the score; it is suitability and explanation.}
\end{frame}

\begin{frame}
\slideTitle{Modern Portfolio Theory: Intuition}
\mainidea{Portfolio construction is about combining assets so the whole portfolio behaves better than isolated holdings.}
\begin{itemize}
  \item Return matters, but risk and correlation also matter.
  \item Diversification works when assets do not move perfectly together.
  \item A portfolio is not a shopping basket; weights determine exposures.
\end{itemize}
\finlens{Robo-advisory turns these ideas into scalable portfolio recommendations under client constraints.}
\end{frame}

\begin{frame}
\slideTitle{Visual Example: Efficient Frontier}
\sourcefigure{figures_pdf/efficient_frontier_cml.pdf}{0.72}{Wikimedia Commons, Capital market line of CAPM, used in Modern Portfolio Theory article.}
\end{frame}


\begin{frame}
\slideTitle{Modern Portfolio Theory: Formula}
\begin{itemize}
  \item For portfolio weights $w$, expected return $\mu$ and covariance $\Sigma$:
\end{itemize}
\[
    \mathbb{E}[r_p]=w^\top\mu, \qquad \sigma_p^2=w^\top\Sigma w.
\]
\begin{itemize}
  \item Mean-variance optimization balances return and risk:
\end{itemize}
\[
\max_w \; w^\top\mu - \frac{\lambda}{2}w^\top\Sigma w.
\]
\techlens{Real implementation needs robust estimates, constraints, trading costs and monitoring. The elegant equation is the start, not the product.}
\end{frame}

\begin{frame}
\slideTitle{Minimum Variance Portfolio}
\begin{itemize}
  \item The minimum variance portfolio minimizes total portfolio variance:
\end{itemize}
\[
\min_w \; w^\top \Sigma w \quad \text{s.t.} \quad \mathbf{1}^\top w=1.
\]
\begin{itemize}
  \item With two risky assets, the optimal weight in asset A is:
\end{itemize}
\[
 w_A^*=\frac{\mathrm{Var}(r_B)-\mathrm{Cov}(r_A,r_B)}{\mathrm{Var}(r_A)+\mathrm{Var}(r_B)-2\mathrm{Cov}(r_A,r_B)}.
\]
\takeaway{Diversification is driven by variance, covariance and correlation.}
\end{frame}
% --------------------------------------------------


\begin{frame}
\slideTitle{Practical Portfolio Constraints}
\begin{itemize}
  \item Long-only or limited shorting.
  \item Asset class, sector, region and issuer concentration limits.
  \item Turnover and transaction-cost constraints.
  \item Liquidity and minimum-trade-size constraints.
  \item Tax and account-specific restrictions.
  \item Compliance and suitability restrictions.
\end{itemize}
\finlens{Constraints are not annoying details. They are how a theoretical portfolio becomes an investable portfolio.}
\end{frame}
% --------------------------------------------------



\begin{frame}
\slideTitle{Rebalancing}
\begin{itemize}
  \item Rebalancing restores a portfolio toward target allocation.
  \item Common rules:
  \begin{itemize}
    \item calendar-based: monthly or quarterly;
    \item threshold-based: rebalance if drift exceeds tolerance;
    \item optimization-based: trade only when benefits exceed costs.
  \end{itemize}
  \item AI can help forecast drift, costs and tax consequences.
\end{itemize}
\takeaway{Rebalancing is a trade-off between discipline, cost, tax and risk control.}
\end{frame}

\begin{frame}
\slideTitle{Lab Demo: Robo Advisor Journey}
\footnotesize
\begin{beamercolorbox}[wd=\linewidth,rounded=true,sep=0.12cm]{visualbox}\scriptsize\textbf{Link: https://portal.aqumon.com/home/} \end{beamercolorbox}
\end{frame}


\begin{frame}
\slideTitle{WealthTech Product Matrix}
\begin{itemize}
  \item Digital onboarding and risk profiling.
  \item Portfolio proposal and suitability check.
  \item Goal-based investing.
  \item Rebalancing and tax-aware trading.
  \item Portfolio analytics, stress testing and performance attribution.
  \item Advisor copilot and client communication.
  \item Compliance monitoring and audit trail.
\end{itemize}
\mainidea{WealthTech is a system of connected services, not a single algorithm.}
\end{frame}

\begin{frame}
\slideTitle{Direct Indexing: Plain-English View}
\mainidea{Direct indexing means owning the individual stocks in an index-like portfolio instead of buying one index fund.}
\begin{itemize}
  \item The investor still aims to track a benchmark.
  \item But individual holdings allow more customization.
  \item Examples:
  \begin{itemize}
    \item exclude restricted securities;
    \item apply ESG or thematic tilts;
    \item control sector or factor exposures;
    \item harvest tax losses at individual security or lot level.
  \end{itemize}
\end{itemize}
\visual{Benchmark index decomposed into individual stock tiles, then personalized exclusions and tilts.}
\end{frame}

\begin{frame}
\slideTitle{Visual Example: Direct Indexing Workflow}
\sourcefigure{figures_pdf/direct_indexing_workflow.pdf}{0.88}{BNY Pershing, \textit{The Democratization of Direct Indexing}, direct-indexing workflow infographic.}
\end{frame}


\begin{frame}
\slideTitle{Direct Indexing as an Optimization Problem}
\begin{itemize}
  \item Let $w$ be portfolio weights and $w_b$ benchmark weights.
  \item Tracking error objective:
\end{itemize}
\[
\min_w \; (w-w_b)^\top\Sigma(w-w_b)
\]
\begin{itemize}
  \item Practical version adds costs and constraints:
\end{itemize}
\[
\min_w \; \mathrm{TE}^2(w)+\lambda_1\mathrm{Turnover}(w)+\lambda_2\mathrm{TaxCost}(w)
\]
\[
\text{s.t. } \mathbf{1}^\top w=1,\quad 0\leq w_i\leq u_i,\quad \text{client constraints}.
\]
\end{frame}

\begin{frame}
\slideTitle{Tax-Loss Harvesting: Plain-English View}
\mainidea{Tax-loss harvesting realizes selected investment losses today to reduce tax, while trying to keep similar market exposure.}
\begin{enumerate}
  \item Identify tax lots with unrealized losses.
  \item Sell selected lots.
  \item Buy replacement exposures to maintain risk profile.
  \item Avoid wash-sale or equivalent restrictions where applicable.
  \item Track realized gains/losses and future cost basis.
\end{enumerate}
\finlens{Tax rules vary by jurisdiction; implementation must be localized and reviewed by qualified tax professionals.}
\end{frame}

\begin{frame}
\slideTitle{Visual Example: Tax-Loss Harvesting Process}
\sourcefigure{figures_pdf/tax_loss_harvesting_process.pdf}{0.88}{Mellon Investments Corporation, \textit{Direct Indexing: The Basics of Tax-Loss Harvesting}, process graphic.}
\end{frame}


\begin{frame}
\slideTitle{Tax Harvesting: Decision Logic}
\begin{itemize}
  \item For tax lot $i$, define unrealized loss:
\end{itemize}
\[
L_i = \max(0, \mathrm{CostBasis}_i - \mathrm{MarketValue}_i).
\]
\begin{itemize}
  \item A trade is attractive only if expected tax benefit exceeds costs and tracking-error impact:
\end{itemize}
\[
\mathrm{NetBenefit}_i \approx \tau L_i - \mathrm{TradingCost}_i - \Delta \mathrm{TE}_i - \mathrm{OperationalPenalty}_i.
\]
\techlens{This is an optimization problem with tax lots, replacement securities, benchmark risk, trading costs and client constraints.}
\end{frame}
% --------------------------------------------------


\begin{frame}
\slideTitle{Direct Indexing + Tax Harvesting + AI}
\begin{itemize}
  \item Direct indexing creates more tax lots and more customization degrees of freedom.
  \item Tax harvesting searches across many lots, securities and replacement baskets.
  \item AI can support:
  \begin{itemize}
    \item personalized constraint extraction from client documents;
    \item opportunity ranking;
    \item replacement-security selection;
    \item explanation of after-tax trade-offs;
    \item monitoring and audit trail.
  \end{itemize}
\end{itemize}
\takeaway{The product value comes from personalization at scale, not from a single clever trade.}
\end{frame}



\begin{frame}
\slideTitle{AI Agents in an Investment Firm}
\mainidea{A real investment process involves many roles. Agent design can mirror this organizational structure.}
\begin{center}
\scriptsize
\begin{tabular}{lll}
\toprule
Role & Main question & Typical output \\
\midrule
Macro analyst & What is the regime? & regime view \\
Quant researcher & Does signal work? & research memo \\
Portfolio manager & How to implement? & target weights \\
Risk manager & What can go wrong? & risk challenge \\
Trader & How to trade? & execution plan \\
Compliance officer & Is it allowed? & approval/checklist \\
Committee secretary & What was decided? & audit record \\
\bottomrule
\end{tabular}
\end{center}
\end{frame}



\begin{frame}
\slideTitle{Agent Role: Macro Analyst}
\begin{itemize}
  \item Monitors macro data, central bank communication, inflation, growth and liquidity conditions.
  \item Uses retrieval over research reports, official data and market indicators.
  \item Produces structured regime views:
\end{itemize}
\[
\text{Regime} \in \{\text{Growth Up},\text{Growth Down}\}\times\{\text{Inflation Up},\text{Inflation Down}\}.
\]
\begin{itemize}
  \item Output should include evidence, confidence and changes from previous view.
\end{itemize}
\end{frame}

\begin{frame}
\slideTitle{Agent Role: Quant Researcher}
\begin{itemize}
  \item Converts investment hypotheses into testable factors.
  \item Pulls data, cleans data and checks survivorship and look-ahead bias.
  \item Runs IC analysis, robustness tests and parameter sweeps.
  \item Produces research memo with code, assumptions and performance diagnostics.
\end{itemize}
\finlens{This is where AI can speed up research, but also where bad automation can create beautiful nonsense.}
\end{frame}

\begin{frame}
\slideTitle{Agent Role: Portfolio Manager}
\begin{itemize}
  \item Translates signals into portfolio weights.
  \item Balances expected return, risk, turnover, liquidity, concentration and client constraints.
  \item Evaluates whether proposed trades are consistent with mandate and market view.
  \item Escalates unusual or high-impact changes to the investment committee.
\end{itemize}
\takeaway{The portfolio manager agent should not only ask ``What has alpha?'' but also ``Can we implement it responsibly?''}
\end{frame}

\begin{frame}
\slideTitle{Agent Role: Risk Manager}
\begin{itemize}
  \item Reviews exposures, drawdown, stress scenarios and tail risk.
  \item Challenges the strategy assumptions.
  \item Checks whether backtest risk is concentrated in a small number of episodes.
  \item Requires kill-switches, monitoring thresholds and post-trade review.
\end{itemize}
\techlens{A strong risk agent should look for hidden beta, crowding, liquidity traps, factor concentration and regime sensitivity.}
\end{frame}

\begin{frame}
\slideTitle{Agent Role: Compliance Officer}
\begin{itemize}
  \item Checks suitability, mandate restrictions and disclosure language.
  \item Reviews whether AI-generated analysis is supported by evidence.
  \item Prevents misleading claims such as unsupported ``AI alpha'' marketing.
  \item Maintains audit trail: prompts, retrieved sources, model outputs, approvals and final actions.
\end{itemize}
\takeaway{Compliance is not the enemy of innovation. It is the seat belt that lets the car go faster safely.}
\end{frame}

\begin{frame}
\slideTitle{Agentic Investment Committee}
\begin{enumerate}
  \item Researcher proposes signal.
  \item Risk agent challenges hidden exposures.
  \item Portfolio agent proposes implementation.
  \item Compliance agent checks mandate and disclosure.
  \item Human committee approves, rejects or requests revision.
\end{enumerate}
\mainidea{Agents do not replace accountability; they structure debate and preserve evidence.}
\takeaway{The system should record dissent and rationale, not only the final output.}
\end{frame}
% --------------------------------------------------


\begin{frame}
\slideTitle{From Simple Investment Idea to Strategy}
\begin{center}
\Large\textit{``Buy companies with improving fundamentals, reasonable valuation and strong cash conversion.''}
\end{center}
\vspace{0.2cm}
\begin{itemize}
  \item Convert plain language into measurable signals:
  \begin{itemize}
    \item earnings growth;
    \item free cash flow growth;
    \item return on invested capital;
    \item valuation multiples;
    \item balance sheet quality.
  \end{itemize}
  \item Then test whether the signal predicts future returns after costs.
\end{itemize}
\end{frame}

\begin{frame}
\slideTitle{Strategy Research Pipeline}
\begin{enumerate}
  \item Define investment universe.
  \item Define signal and data sources.
  \item Clean data and avoid biases.
  \item Analyze factor IC and monotonicity.
  \item Tune parameters and rebalance frequency.
  \item Construct portfolio weights.
  \item Backtest with realistic costs and constraints.
  \item Attribute performance and risk.
  \item Prepare deployment and monitoring.
\end{enumerate}
\visual{Pipeline from idea to data, factor test, optimization, backtest, attribution and deployment.}
\end{frame}




\begin{frame}
\slideTitle{AI Agent for Investment Research Implementation}
\begin{itemize}
  \item User gives simple idea in natural language.
  \item Agent converts idea into structured research plan.
  \item Data agent retrieves and validates data.
  \item Quant agent creates factors and runs IC analysis.
  \item Optimization agent builds portfolios under constraints.
  \item Backtest agent evaluates costs, risk and performance.
  \item Attribution agent explains return drivers.
  \item Compliance agent checks disclosures and suitability.
\end{itemize}
\end{frame}





\begin{frame}
\slideTitle{Summary}
\begin{itemize}
  \item AI evolved from rule-based systems to machine learning, deep learning, foundation models, agents and world models.
  \item FinTech applications require prediction, generation, reasoning, planning and governance.
  \item Robo-advisory, direct indexing and tax harvesting are examples of algorithmic personalization at scale.
  \item Agentic workflows can mimic the decision process of an investment firm, but must preserve accountability.
  \item A sound investment AI workflow goes from idea to data, IC, tuning, weighting, backtesting, attribution and monitoring.
\end{itemize}
\end{frame}

\begin{frame}
\slideTitle{Questions?}
\centering
\includegraphics[width=0.36\textwidth]{figures_pdf/wechat_group.pdf}
\end{frame}




\begin{frame}
\slideTitle{References and Further Reading}
\scriptsize
\begin{itemize}
  \item Turing, A. M. (1950), ``Computing Machinery and Intelligence.''
  \item Vaswani et al. (2017), ``Attention Is All You Need.''
  \item Sutton and Barto, ``Reinforcement Learning: An Introduction.''
  \item Murphy, ``Machine Learning: A Probabilistic Perspective.''
  \item Goodfellow, Bengio and Courville, ``Deep Learning.''
  \item Research and industry materials on foundation models, RAG, AI agents and world models.
  \item World Bank, ``Robo-Advisors: Investing through Machines.''
  \item SEC, robo-adviser guidance and fiduciary duty materials.
  \item Materials from major asset managers on direct indexing and tax-loss harvesting.
\end{itemize}
\end{frame}



\end{document}
